\documentclass[12pt]{article}
\usepackage{requerimientos}
%\usepackage[utf8x]{inputenc}
\usepackage{booktabs}
\usepackage[spanish]{babel}

%\usepackage{amssymb,amsmath,amsthm,amsfonts}
\usepackage{calc}
\usepackage{graphicx}
\usepackage{subfigure}
\usepackage{gensymb}
\usepackage{natbib}
\usepackage{url}
\usepackage[utf8x]{inputenc}
\usepackage{amsmath}
\usepackage{graphicx}
\graphicspath{{images/}}
\usepackage{parskip}
\usepackage{fancyhdr}
\usepackage{vmargin}
\setcounter{secnumdepth}{4}

\setmarginsrb{3 cm}{2.5 cm}{3 cm}{2.5 cm}{1 cm}{1.5 cm}{1 cm}{1.5 cm}

\title{Requerimientos y casos de uso}					% Titulo
\author{Ingeniería de Software}					% Autor


\makeatletter
\let\thetitle\@title
\let\theauthor\@author
\let\thedate\@date
\makeatother

\pagestyle{fancy}
\fancyhf{}
\rhead{\theauthor}
\lhead{\thetitle}
\cfoot{\thepage}

\begin{document}

%%%%%%%%%%%%%%%%%%%%%%%%%%%%%%%%%%%%%%%%%%%%%%%%%%%%%%%%%%%%%%%%%%%%%%%%%%%%%%%%%%%%%%%%%

\begin{titlepage}
	\centering
    \vspace*{0.0 cm}
    \includegraphics[scale = 0.5]{fotoEscudoESCOM.jpg}\\[.0 cm]	% Logo Universidad
    \textsc{\LARGE Escuela Superior de Cómputo}\\[2.0 cm]	% Nombre Universidad
	\textsc{\Large Ingeniería De Software}\\[0.5 cm]				% Codigo Curso
	\textsc{\large  MALDONADO CASTILLO M. EN C. IDALIA}\\[0.5 cm]		% Nombre Curso
	\rule{\linewidth}{0.2 mm} \\[0.4 cm]
	{ \huge \bfseries \thetitle}\\
	\rule{\linewidth}{0.2 mm} \\[1.5 cm]
	
	\begin{minipage}{0.4\textwidth}
		\begin{center} \large
			\emph{Integrantes de equipo:}\\
            \emph{Castro Flores Marcela}
	    \emph{Chavez Chavez Javier}
            \emph{Guzmán Chávez Juan Antonio}
            \emph{Maldonado Romo Alberto}
            \emph{Murga Dionicio Rubén}
            \emph{}
			\theauthor\linebreak
			\end{center}
	\end{minipage}\\[2 cm] 
	\vfill
	
\end{titlepage}

%%%%%%%%%%%%%%%%%%%%%%%%%%%%%%%%%%%%%%%%%%%%%%%%%%%%%%%%%%%%%%%%%%%%%%%%%%%%%%%%%%%%%%%%%

\tableofcontents
\pagebreak

%%%%%%%%%%%%%%%%%%%%%%%%%%%%%%%%%%%%%%%%%%%%%%%%%%%%%%%%%%%%%%%%%%%%%%%%%%%%%%%%%%%%%%%%%
\section{Requerimientos Funcionales}

\subsection{Requerimientos globales}
\begin{FRequirements}
\FRitem{RFG1}{Bitácora de transacciones}{El sistema mantendrá una bitácora de todas las transacciones de trámites finalizadas a lo largo de un semestre, así hayan sido aceptadas o rechazadas.
}
\FRitem{RFG2}{Importación y exportación de archivos}{El sistema permitirá exportar e importar datos que se encuentran en archivos de tipo .xsl .
}
\end{FRequirements}

\subsection{Jefa de Gestión Escolar}
\begin{FRequirements}
\FRitem{RFJGE1}{Reporte global de transacciones}{El sistema permitirá la visualización de un reporte global de las transacciones de trámites hechas diariamente, mensualmente y semestralmente.
}
\FRitem{RFJGE2}{Control de permisos}{El sistema permitirá al jefe de gestión escolar delegar o autorizar permisos a un analista.
}
\FRitem{RFJGE3}{Trámites realizados}{El sistema permitirá la visualización gráfica de los trámites tanto rechazados,aceptados, recogidos, no recogidos así como por fecha y alumno.
}
\end{FRequirements}


\subsection{Autoridad Pertinente}
\begin{FRequirements}
\FRitem{RFA1}{Reporte global de transacciones}{
El sistema permitirá la visualización de un reporte global de las transacciones hechas diariamente, mensualmente y semestralmente.
}
\end{FRequirements}


\subsection{Analista}
\begin{FRequirements}
\FRitem{RFA1}{Actualización del estado de trámite}{
El sistema permitirá al actor actualizar el estado en que se encuentra los trámites manualmente.
}
\FRitem{RFA2}{Control de trámites solicitados}{
El sistema permitirá controlar el número de trámites solicitados mediante un mecanismo de aceptar o rechazar peticiones.
}
\FRitem{RFA3}{Recuperación de contraseña}{
El sistema permitirá la recuperación de contraseña olvidada.
}
\FRitem{RFA4}{Autorización de trámite}{
El sistema almacenará la solicitud de boleta o constancia hecha por el actor hasta ser aceptada.
}
\end{FRequirements}


\subsection{Alumno / Egresado}
\begin{FRequirements}
\FRitem{RFAE1}{Solicitud de trámite}{
El sistema permitirá al actor solicitar uno o más trámites.
}
\FRitem{RFAE2}{Consulta estado de trámite}{
El sistema permitirá consultar el estado en que se encuentra el trámite.
}
\end{FRequirements}

\subsection{Departamentos escolares}
\begin{FRequirements}
\FRitem{RFDE1}{Solicitud de constancias}{
El sistema permitirá al personal de otro departamento hacer solicitudes de constancias vía memorándum.
}
\FRitem{RFDE2}{Recuperación de contraseña}{
El sistema permitirá la recuperación de contraseña olvidada.
}
\end{FRequirements}

\section{Requerimientos no funcionales}
\begin{NFRequieriments}
\NFRitem{RNF1}{Seguridad}{
\begin{itemize}
\item El sistema utilizará un captcha para evitar inserciones ficticias.
\end{itemize}
}
\end{NFRequieriments}

\begin{NFRequieriments}
\NFRitem{2}{Usabilidad}{
\begin{itemize}
\item El sistema tendrá un diseño  de la interfaz con la gama de colores y logos predefinidos por el IPN y ESCOM.
\item El sistema será hecho en plataforma web
\item El sistema será compatible con los navegadores Chrome.
\item El sistema tendrá un diseño responsivo,es decir estará disponible para computadoras, tablets y smartphones.
\item El sistema se apegará a los lineamientos de las 8 reglas de Shneiderman.
\end{itemize}
}
\end{NFRequieriments}

\begin{NFRequieriments}
\NFRitem{3}{Disponibilidad}{
\begin{itemize}
\item El sistema podrá ser accesado en cualquier momento, soportando la concurrencia que dependerá del servidor proporcionado.
\end{itemize}
}
\end{NFRequieriments}

\begin{NFRequieriments}
\NFRitem{4}{Mantenibilidad}{
\begin{itemize}
\item  El sistema debe proporcionar alertas que indiquen al usuario errores cometidos o información de una ejecución correcta. 
\item El sistema contará con un apartado con información de contacto de los desarrolladores para aclaración.
\end{itemize}
}
\end{NFRequieriments}

\begin{NFRequieriments}
\NFRitem{5}{Extensibilidad}{
\begin{itemize}
\item El sistema podrá tener un crecimiento en el futuro debido a que su desarrollo será por módulos.
\end{itemize}
}
\end{NFRequieriments}

\end{document}
